\clearpage
\section*{摘要}
\addcontentsline{toc}{section}{摘要}
在数据驱动的社会背景下，数据非法获取、篡改或泄露的风险显著增加，不仅威胁个人财产安全，还可能影响社会稳定和国家安全。
传统加密技术虽然能够保护数据传输内容，却难以隐藏通信行为本身。相比之下，图像隐写技术虽然能够隐藏秘密信息和通信行为的存在性，但部分基于深度学习的方案依赖专用模型训练或额外组件，增加了部署成本。
针对这些问题，受深度神经网络（Deep Neural Network，DNN）和密码学融合启发，本文提出一种基于图像密钥控制DNN状态的安全隐蔽信息表示框架（Secure and Stealthy Information Representation，SSIR）。

该框架的发送方和接收方预先共享密钥。发送方首先使用分组编码将消息标识为现有DNN中选定神经的目标状态，随后结合共享密钥，通过对抗优化生成能够诱导目标状态的密文图像。接收方将密文图像与共享密钥融合后输入相同的DNN，通过解码神经元状态恢复原始消息。
由于密文图像可能存在直接视觉泄露和间接DNN状态泄露风险，方案设计了联合损失优化机制，以抑制图像中的可检测信息痕迹。进一步，我们从理论上分析了SSIR的安全性。针对查询有界攻击者，分别推到其IND-CPA与IND-CCA1优势界，
并给出有界查询安全所需的分布条件与每密钥查询次数限制。

与依赖专用模型训练的部分现有方法相比，SSIR无需修改DNN参数或进行额外模型训练，即可利用现有模型完成信息表示与恢复，并保留其原有功能。实验基于ImageNet数据集和真实图像，从消息恢复可靠性、信息泄漏抑制、攻击鲁棒性、运行效率和模型兼容性等方面对方案进行评估。
实验结果表明，分组编码将所测配置下的生成成功率由79\%提高至100\%;通过调整优化参数，直接视觉泄漏代理指标最高降低99.79\%，启用间接泄漏约束后，神经元状态偏差平均降低56.02\%;对于150位消息，密文图像生成时间约为300～500 ms，消息恢复时间约为10 ms。
SSIR能够适配多种常见DNN，并且保持模型原有功能。通过分析不同分组大小和学习率对生成成功率，图像质量及生成时间的影响，为实际应用中的参数选择提供了参考。

\vspace{\baselineskip}
\noindent\textbf{关键词：}深度神经网络；信息安全；隐蔽信息传输；对抗样本